\documentclass[12pt,a4paper]{article}
\usepackage[utf8]{inputenc}
\usepackage[spanish,german]{babel}
\usepackage[T1]{fontenc}
\usepackage{geometry}
\geometry{a4paper, margin=2.5cm}
\usepackage{titlesec}
\usepackage{enumitem}
\usepackage{multicol}
\usepackage{csquotes}

\titleformat{\section}
{\Large\bfseries}
{}
{0em}
{}[\titlerule]

\titleformat{\subsection}
{\large\bfseries}
{}
{0em}
{}

\setlength{\parindent}{0pt}
\setlength{\parskip}{0.5em}

\begin{document}

\title{\textbf{DIALOG: Decir la Hora}}
\author{Spanisch A1/A2}
\date{}
\maketitle

\textbf{Tema:} Decir la hora \\
\textbf{Nivel:} A1 \\
\textbf{Complejidad:} Einfach \\
\textbf{Palabras:} 723

\vspace{1em}

\textit{Nota: Se recomienda revisar primero el vocabulario antes de leer el diálogo.}

\vspace{1em}

\section*{Diálogo}

\textbf{Ana:} ¡Hola, Diego! ¿Qué hora es? No tengo reloj.

\textbf{Diego:} Hola, Ana. Déjame ver... Son las tres y cuarto.

\textbf{Ana:} ¿Las tres y cuarto? ¿Estás seguro? Creo que son las tres y media.

\textbf{Diego:} Tienes razón, perdón. Son las tres y media. Mi reloj está un poco atrasado.

\textbf{Ana:} No te preocupes. ¿A qué hora empieza la clase de español?

\textbf{Diego:} La clase empieza a las cuatro en punto. Tenemos media hora todavía.

\textbf{Ana:} Perfecto. Entonces tenemos tiempo para tomar un café. ¿A qué hora quieres quedar?

\textbf{Diego:} Podemos quedar a las tres y cuarenta y cinco. Así llegamos a tiempo a la clase.

\textbf{Ana:} De acuerdo. Oye, ¿me puedes decir qué hora es ahora?

\textbf{Diego:} Claro, son las tres y treinta y cinco. Faltan diez minutos para las tres y cuarenta y cinco.

\textbf{Ana:} Gracias. ¿A qué hora termina la clase de español?

\textbf{Diego:} La clase termina a las cinco y media. Después tengo otra clase a las seis menos cuarto.

\textbf{Ana:} ¿A las seis menos cuarto? Eso es a las cinco y cuarenta y cinco, ¿verdad?

\textbf{Diego:} Exacto. En español también puedes decir \enquote{las cinco y cuarenta y cinco} o \enquote{las seis menos cuarto}. Ambas formas son correctas.

\textbf{Ana:} Interesante. ¿Y cómo se dice cuando son las doce del mediodía?

\textbf{Diego:} Se dice \enquote{es mediodía} o \enquote{son las doce del mediodía}. Y a las doce de la noche se dice \enquote{es medianoche} o \enquote{son las doce de la noche}.

\textbf{Ana:} Entiendo. ¿Me puedes decir qué hora es ahora?

\textbf{Diego:} Son las tres y cuarenta. Ya es hora de ir al café.

\textbf{Ana:} Tienes razón. Vamos. Oye, ¿a qué hora abren los cafés aquí?

\textbf{Diego:} La mayoría de los cafés abren a las siete de la mañana y cierran a las diez de la noche.

\textbf{Ana:} ¿Y los fines de semana?

\textbf{Diego:} Los sábados y domingos abren más tarde, normalmente a las ocho de la mañana.

\textbf{Ana:} ¿A qué hora desayunas normalmente?

\textbf{Diego:} Desayuno a las ocho y media de la mañana. ¿Y tú?

\textbf{Ana:} Yo desayuno más temprano, a las siete y cuarto. Me levanto muy temprano.

\textbf{Diego:} ¿A qué hora te levantas?

\textbf{Ana:} Me levanto a las seis y media de la mañana. ¿Y tú?

\textbf{Diego:} Yo me levanto a las siete en punto. No me gusta levantarme muy temprano.

\textbf{Ana:} Entiendo. ¿A qué hora almuerzas?

\textbf{Diego:} Almuerzo a la una y media. ¿Y tú?

\textbf{Ana:} Yo almuerzo más tarde, a las dos y cuarto. Me gusta almorzar tranquila.

\textbf{Diego:} ¿Y a qué hora cenas?

\textbf{Ana:} Ceno a las ocho de la noche. ¿Y tú?

\textbf{Diego:} Yo ceno más tarde, a las nueve y media. En mi país cenamos más tarde que en otros países.

\textbf{Ana:} Sí, lo he notado. Oye, ¿qué hora es ahora? No quiero llegar tarde a la clase.

\textbf{Diego:} Son las tres y cincuenta. Faltan diez minutos para las cuatro. Tenemos que irnos ya.

\textbf{Ana:} Tienes razón. Vamos rápido. ¿Cuánto tiempo tardamos en llegar?

\textbf{Diego:} Tardamos unos cinco minutos caminando. Llegamos justo a tiempo.

\textbf{Ana:} Perfecto. Gracias por decirme la hora todo el tiempo.

\textbf{Diego:} De nada. Es importante llegar puntual a las clases.

\textbf{Ana:} Sí, tienes razón. La puntualidad es muy importante.

\newpage

\section*{Vocabulario}

\begin{multicols}{2}
\begin{itemize}[leftmargin=*]
    \item \textbf{reloj} - Uhr
    \item \textbf{¿Qué hora es?} - Wie spät ist es?
    \item \textbf{son las [hora]} - es ist [Uhr] (für 2-12)
    \item \textbf{es la [hora]} - es ist [Uhr] (für 1)
    \item \textbf{y cuarto} - Viertel nach
    \item \textbf{y media} - halb (30 Minuten nach)
    \item \textbf{en punto} - genau, Punkt
    \item \textbf{atrasado/a} - verspätet, nachgehend
    \item \textbf{empezar} - beginnen
    \item \textbf{terminar} - enden, beenden
    \item \textbf{menos cuarto} - Viertel vor
    \item \textbf{mediodía} - Mittag
    \item \textbf{medianoche} - Mitternacht
    \item \textbf{mañana} - Morgen
    \item \textbf{tarde} - Nachmittag, später
    \item \textbf{noche} - Nacht, Abend
    \item \textbf{abrir} - öffnen
    \item \textbf{cerrar} - schließen
    \item \textbf{desayunar} - frühstücken
    \item \textbf{almorzar} - zu Mittag essen
    \item \textbf{cenar} - zu Abend essen
    \item \textbf{levantarse} - aufstehen
    \item \textbf{faltar} - fehlen, noch sein bis
    \item \textbf{tardar} - dauern, brauchen (Zeit)
    \item \textbf{puntual} - pünktlich
    \item \textbf{puntualidad} - Pünktlichkeit
    \item \textbf{quedar} - sich treffen, verabreden
    \item \textbf{seguro/a} - sicher
    \item \textbf{preocuparse} - sich sorgen
\end{itemize}
\end{multicols}

\newpage

\section*{Preguntas de Comprensión}

\begin{enumerate}[leftmargin=*]
    \item ¿Qué hora es cuando Ana pregunta por primera vez?
    \item ¿A qué hora empieza la clase de español?
    \item ¿A qué hora quieren quedar Ana y Diego para tomar un café?
    \item ¿A qué hora termina la clase de español?
    \item ¿A qué hora tiene Diego otra clase después de español?
    \item ¿Cómo se dice \enquote{las seis menos cuarto} de otra forma?
    \item ¿Cómo se dice \enquote{las doce del mediodía}?
    \item ¿A qué hora abren normalmente los cafés?
    \item ¿A qué hora cierran los cafés normalmente?
    \item ¿A qué hora abren los cafés los fines de semana?
    \item ¿A qué hora desayuna Diego?
    \item ¿A qué hora se levanta Ana?
    \item ¿A qué hora almuerza Ana?
    \item ¿A qué hora cena Diego?
    \item ¿Por qué Diego dice que en su país cenan más tarde?
    \item ¿Cuánto tiempo tardan Ana y Diego en llegar a la clase?
    \item ¿Qué dice Diego sobre la puntualidad?
    \item ¿Por qué Ana pregunta varias veces qué hora es?
    \item ¿Qué problema tiene el reloj de Diego?
    \item ¿Qué hacen Ana y Diego antes de ir a la clase?
\end{enumerate}

\newpage

\section*{Verbo Irregular: DECIR (sagen)}

\textit{Nota: \enquote{vosotros/as} se usa solo en España, no en español sudamericano. Se incluye aquí para completitud.}

\begin{multicols}{2}

\subsection*{Presente}
\begin{itemize}[leftmargin=*]
    \item yo - digo
    \item tú - dices
    \item él/ella/usted - dice
    \item nosotros/as - decimos
    \item vosotros/as - decís
    \item ustedes - dicen
    \item ellos/ellas - dicen
\end{itemize}

\subsection*{Pretérito Indefinido}
\begin{itemize}[leftmargin=*]
    \item yo - dije
    \item tú - dijiste
    \item él/ella/usted - dijo
    \item nosotros/as - dijimos
    \item vosotros/as - dijisteis
    \item ustedes - dijeron
    \item ellos/ellas - dijeron
\end{itemize}

\subsection*{Pretérito Imperfecto}
\begin{itemize}[leftmargin=*]
    \item yo - decía
    \item tú - decías
    \item él/ella/usted - decía
    \item nosotros/as - decíamos
    \item vosotros/as - decíais
    \item ustedes - decían
    \item ellos/ellas - decían
\end{itemize}

\columnbreak

\subsection*{Futuro Simple}
\begin{itemize}[leftmargin=*]
    \item yo - diré
    \item tú - dirás
    \item él/ella/usted - dirá
    \item nosotros/as - diremos
    \item vosotros/as - diréis
    \item ustedes - dirán
    \item ellos/ellas - dirán
\end{itemize}

\subsection*{Pretérito Perfecto}
\begin{itemize}[leftmargin=*]
    \item yo - he dicho
    \item tú - has dicho
    \item él/ella/usted - ha dicho
    \item nosotros/as - hemos dicho
    \item vosotros/as - habéis dicho
    \item ustedes - han dicho
    \item ellos/ellas - han dicho
\end{itemize}

\end{multicols}

\newpage

\section*{Explicación Gramatical}

\textbf{Die Uhrzeit auf Spanisch -- Cómo decir la hora}

Im Spanischen gibt es verschiedene Möglichkeiten, die Uhrzeit auszudrücken. Die wichtigsten Regeln sind:

\textbf{1. Grundstruktur:}
\begin{itemize}[leftmargin=*]
    \item Für 1 Uhr: \textbf{Es la una} (Es ist ein Uhr)
    \item Für 2-12 Uhr: \textbf{Son las [Zahl]} (Es ist [Zahl] Uhr)
    \item Beispiel: \enquote{Son las tres} (Es ist drei Uhr)
\end{itemize}

\textbf{2. Minuten hinzufügen:}
\begin{itemize}[leftmargin=*]
    \item \textbf{y [Minuten]} - nach (z.B. \enquote{son las tres y cinco} = drei Uhr fünf)
    \item \textbf{y cuarto} - Viertel nach (z.B. \enquote{son las tres y cuarto} = drei Uhr fünfzehn)
    \item \textbf{y media} - halb (z.B. \enquote{son las tres y media} = drei Uhr dreißig)
    \item \textbf{menos [Minuten]} - vor (z.B. \enquote{son las cuatro menos cinco} = fünf vor vier)
    \item \textbf{menos cuarto} - Viertel vor (z.B. \enquote{son las cuatro menos cuarto} = Viertel vor vier)
\end{itemize}

\textbf{3. Tageszeiten:}
\begin{itemize}[leftmargin=*]
    \item \textbf{de la mañana} - morgens (bis 12 Uhr)
    \item \textbf{del mediodía} - mittags (12 Uhr)
    \item \textbf{de la tarde} - nachmittags (13-19 Uhr)
    \item \textbf{de la noche} - abends/nachts (20-24 Uhr)
    \item \textbf{medianoche} - Mitternacht (24 Uhr)
\end{itemize}

\textbf{4. Besondere Ausdrücke:}
\begin{itemize}[leftmargin=*]
    \item \textbf{en punto} - genau, Punkt (z.B. \enquote{son las cuatro en punto})
    \item \textbf{mediodía} - Mittag (12 Uhr mittags)
    \item \textbf{medianoche} - Mitternacht (12 Uhr nachts)
\end{itemize}

\textbf{Ejemplos del texto:}
\begin{itemize}[leftmargin=*]
    \item \enquote{Son las tres y cuarto} -- Es ist Viertel nach drei
    \item \enquote{Son las tres y media} -- Es ist halb vier
    \item \enquote{La clase empieza a las cuatro en punto} -- Die Klasse beginnt um vier Uhr genau
    \item \enquote{Las seis menos cuarto} -- Viertel vor sechs (oder: \enquote{las cinco y cuarenta y cinco})
    \item \enquote{Es mediodía} -- Es ist Mittag
    \item \enquote{Me levanto a las seis y media de la mañana} -- Ich stehe um halb sieben morgens auf
    \item \enquote{Ceno a las ocho de la noche} -- Ich esse um acht Uhr abends zu Abend
\end{itemize}

\textbf{Wichtiger Hinweis:} In südamerikanischem Spanisch werden beide Formen (\enquote{menos cuarto} und die direkte Angabe wie \enquote{cinco y cuarenta y cinco}) verwendet, aber die direkte Form ist häufiger.

\end{document}
